\section{Introduction}
Understanding and predicting how the world evolves from observations is one of the central goals of artificial intelligence \cite{schmidhuber2015learning,ha2018world,kim2020active}.
Both dynamical systems theory \cite{bertsekas2012dynamic,hafner2023mastering} and reinforcement learning \cite{sutton2018reinforcement} typically model the world as a latent-state dynamical process, where the environment evolves through state transitions driven by actions. From this perspective, visual observations are merely partial and noisy projections of the true system state.
Therefore, learning a predictive model of the world requires inferring latent states and modeling their action-conditioned state transitions.
Such world models are crucial for enabling agents to plan, reason, and interact with complex environments over long horizons.


Recent years have witnessed significant progress in video generation and world models \cite{wan2025wan,hacohen2026ltx,team2026advancing}.
Many recent approaches~\cite{matrixgame2,ji2025memflow,yume1_5} attempt to learn environment dynamics from large-scale video datasets by training generative models that predict future frames conditioned on past observations and actions.
However, despite the increasing capability of such models, existing datasets remain insufficient for effectively learning structured action-conditioned dynamics.
Most existing datasets provide only simple action annotations with limited semantic meaning, such as basic movements or camera rotations \cite{sekai,spatialvid}.
Moreover, the effects of these actions are often directly observable in the visual observations.
For example, the action “move left” is typically reflected in the video as a corresponding change in viewpoint.

However, in many cases, actions are not defined through explicit observation variations but instead manifest through implicit state transitions.
For instance, the action ``shoot'' implicitly affects internal state variables such as the ``remaining ammunition count''.
This state cannot be reliably inferred from visual observations alone, yet it plays a crucial role in determining future visual outcomes.
When the remaining ammunition reaches zero, executing the shoot action will no longer produce firing effects or projectiles, leading to visual results that differ significantly from those observed when ammunition is available.
Such coupling makes it difficult for models to disentangle state transitions from observation variations, thereby hindering the learning of stable and interpretable world dynamics.
As a result, current models often perform poorly in long-horizon prediction tasks, where small errors accumulate over time and eventually lead to noticeable inconsistencies or instability in the generated results.

In this paper, we propose WildWorld, a large-scale video dataset for action-conditioned world modeling with explicit state annotations.
The dataset is automatically collected from the photorealistic AAA action role-playing game \textit{Monster Hunter: Wilds}. WildWorld features a rich and semantically meaningful action space containing over 450 actions, including movement, attacks, and skill casting.
To facilitate data collection, we develop a bespoke toolchain capable of recording per-frame ground-truth annotations, including player actions, character skeletons, world states, camera poses, depth maps, etc.
The toolchain is integrated with an automated gameplay pipeline, allowing the dataset to scale easily to over 108M frames of gameplay footage while covering diverse interactive scenarios.
By capturing complex interactions and the underlying state transitions, WildWorld enables the study of long-horizon compositional action sequences and their effects on evolving world states, providing a valuable foundation for building, training, and systematically evaluating state-aware interactive world models.

Furthermore, we derive WildBench, a benchmark constructed from WildWorld for evaluating interactive world models. WildBench introduces two key evaluation metrics: Action Following and State Alignment.
Specifically, Action Following measures the agreement between generated videos and the ground-truth sub-actions.%, which is automatically evaluated using large vision-language models.
State alignment quantitatively measures the accuracy of state transitions by tracking skeletal keypoints in the generated videos and comparing them with the corresponding ground-truth annotations.
We design several baseline models for state-aware video generation and compare them with existing approaches on WildBench.
The experimental results reveal the limitations of current models and provide insights for future research, particularly in improving state transition modeling and long-horizon consistency.

To summarize, our contributions are threefold:
(1) We propose \textbf{WildWorld}, a large-scale video dataset comprising over 108M frames, with a rich action space and diverse frame-level ground truth annotations, including player actions, character skeletons, world states, camera poses, depth, etc.
(2) We curate \textbf{WildBench}, a benchmark for evaluating interactive world models, featuring two carefully designed metrics: Action Following and State Alignment.
(3) We conduct extensive experiments and analysis on WildBench, which provide insights into the future development of interactive world models.
